\section{Experiments}
\label{sec:experiments}

We design experiments to answer twelve research questions:
\begin{description}[leftmargin=*,topsep=2pt,itemsep=2pt]
  \item[\textbf{Q1}] Does UCDPA improve accuracy over COME and other TTA baselines? (Table~\ref{tab:main})
  \item[\textbf{Q2}] Does UCDPA improve calibration? (Table~\ref{tab:calibration})
  \item[\textbf{Q3}] Does UCDPA improve error detection? (Table~\ref{tab:auroc})
  \item[\textbf{Q4}] How does UCDPA perform across severity levels? (Table~\ref{tab:severity})
  \item[\textbf{Q5}] How does UCDPA perform across corruption families? (Table~\ref{tab:family})
  \item[\textbf{Q6}] How much does each component contribute? (Table~\ref{tab:ablation})
  \item[\textbf{Q7}] Is the UCB score well-separated between correct and incorrect predictions? (\S\ref{sec:exp:ucb-analysis})
  \item[\textbf{Q8}] Does UCDPA generalize across backbones (CIFAR-ResNet-18) and datasets (CIFAR-100-C)? (Table~\ref{tab:cross-backbone})
  \item[\textbf{Q9}] How sensitive is UCDPA to the routing thresholds $\tau_c, \tau_u$ and temperatures $T_c, T_u$? (Table~\ref{tab:sensitivity})
  \item[\textbf{Q10}] How does UCDPA perform under batch size 1 and small batch sizes? (Table~\ref{tab:batch-size})
  \item[\textbf{Q11}] Does scale normalization matter? (Table~\ref{tab:scale-norm})
  \item[\textbf{Q12}] What is the runtime and memory overhead relative to COME? (\S\ref{sec:exp:efficiency})
\end{description}

\subsection{Experimental Setup}
\label{sec:exp:setup}

\paragraph{Dataset.}
We evaluate on ImageNet-C~\citep{hendrycks2019robustness}, the standard TTA benchmark. ImageNet-C contains 19 corruption types (each applied to the 50{,}000-image ImageNet validation set) across 5 severity levels, grouped into 4 families: \emph{noise} (Gaussian, shot, impulse, speckle), \emph{blur} (defocus, glass, motion, zoom, frost), \emph{weather} (snow, frost, fog, brightness), and \emph{digital} (contrast, elastic, pixelate, JPEG), plus an extra set. We follow the standard TTA protocol: the model is adapted online on the full 50{,}000-image stream for each corruption/severity, with the stream order fixed across methods. We report mean corruption accuracy averaged over all 19 corruptions at severity 1--5, and per-family and per-severity breakdowns.

\paragraph{Backbone.}
We use ViT-B/16~\citep{dosovitskiy2021vit} pretrained on ImageNet-1k, following the standard TTA setup~\citep{niu2023sar,come2024come}. We adapt only the LayerNorm affine parameters (as in~\citep{niu2023sar,come2024come}), leaving all other parameters frozen. The penultimate feature dimension is $d=768$.

\paragraph{Baselines.}
We compare against seven methods:
\begin{enumerate}[leftmargin=*,topsep=2pt,itemsep=2pt]
  \item \textbf{Source-only}: no adaptation; the source-pretrained model is evaluated directly.
  \item \textbf{Tent}~\citep{wang2021tent}: entropy minimization on LayerNorm parameters.
  \item \textbf{EATA}~\citep{niu2022eata}: reliable sample-weighted entropy plus Fisher regularization.
  \item \textbf{SAR}~\citep{niu2023sar}: sharpness-aware entropy minimization.
  \item \textbf{COME}~\citep{come2024come}: conservative opinion-entropy minimization (the primary baseline).
  \item \textbf{Tent-COME}: Tent augmented with the COME evidential opinion entropy (a hybrid baseline isolating the contribution of the conservative entropy beyond Tent).
  \item \textbf{UCDPA} (ours): the full method of Algorithm~\ref{alg:ucdpa}.
\end{enumerate}

\paragraph{Hyperparameters.}
For UCDPA we use $\tau_c=0.7$, $\tau_u=0.3$, $T_c=T_u=0.1$, $\lambda_\mathrm{come}=0.03$, $\lambda_\mathrm{cal}=0.02$, $\lambda_\mathrm{fpa}=0$ (disabled), prototype momentum $m=0.9$, EMA class-counter momentum $\beta=0.99$, FPA temperature $T_\mathrm{fpa}=0.1$, Adam learning rate $2\times10^{-4}$, and batch size $B=128$. All baselines use the hyperparameters reported in their original papers (tuned for ViT-B/16 on ImageNet-C). Full hyperparameters are in Appendix~\ref{app:hyperparameters}.

\paragraph{Protocol.}
We run each method with 3 random seeds (controlling batch ordering and any stochastic augmentation). We report the mean. We compute 95\% bootstrap confidence intervals (10{,}000 resamples) for the headline comparison and a paired permutation test for statistical significance (Section~\ref{sec:exp:main}). All methods see the same stream order per seed.

\paragraph{Metrics.}
We report: (i) \textbf{mean corruption accuracy} (mean over 19 corruptions, severity 1--5); (ii) \textbf{Expected Calibration Error (ECE)} and \textbf{Maximum Calibration Error (MCE)}~\citep{guo2017calibration} with 15 bins; (iii) \textbf{Predictive Entropy} $\E[H(p)]$ and \textbf{Mutual Information} $\E[H(p)] - \E_{\alpha}[H(p)]$ between the Dirichlet-mean categorical and the Dirichlet-marginal categorical~\citep{smith2018understanding}; (iv) \textbf{AUROC} for error detection, where the binary task is correct-vs-incorrect (using ground-truth labels only for evaluation, never for adaptation) and the score is $1 - u_i$ (or $c_i$ for non-evidential baselines).

\subsection{Q1: Main Accuracy Comparison}
\label{sec:exp:main}

\begin{table}[t]
\centering
\caption{Main accuracy comparison on ImageNet-C (ViT-B/16, mean over 19 corruptions, severity 1--5, 3 seeds). Higher is better. $\Delta = \text{UCDPA} - \text{COME}$ is the accuracy gain in percentage points. $p$-value from a paired permutation test (10{,}000 permutations) of the null hypothesis $H_0\colon \Delta = 0$; 95\% CI from 10{,}000 nonparametric bootstrap resamples over the 57 (corruption, seed) units, reported on $\Delta$ (in pp).}
\label{tab:main}
\begin{tabular}{lcc}
\toprule
Method & Accuracy (\%) & Notes \\
\midrule
Source-only & \sourceacc & No adaptation \\
Tent~\citep{wang2021tent} & \tentacc & Entropy minimization \\
EATA~\citep{niu2022eata} & \eataacc & Reliable entropy + Fisher \\
SAR~\citep{niu2023sar} & \saracc & Sharpness-aware entropy \\
COME~\citep{come2024come} & \comeacc & Conservative opinion entropy \\
Tent-COME & \tentcomeacc & Tent + opinion entropy \\
\textbf{UCDPA (ours)} & \ucdpaacc & Dual-path + UCB + cal + FPA \\
\midrule
$\Delta$ (UCDPA $-$ COME) & \ucdpadelta & $p < \pvalue$ \\
95\% CI of $\Delta$ (pp) & $[\ucdpacilow, \ucdpacihigh]$ & Bootstrap \\
\bottomrule
\end{tabular}
\end{table}


Table~\ref{tab:main} reports mean corruption accuracy on ImageNet-C with ViT-B/16. The key findings are:

\textbf{UCDPA improves over COME by \ucdpadelta{} percentage points} (\ucdpaacc{} vs.\ \comeacc{}). This improvement is \textbf{statistically significant}: a paired permutation test (10{,}000 permutations, paired over (corruption, seed)) of $H_0\colon \Delta = 0$ yields $p < \pvalue{}$, and the 95\% bootstrap confidence interval for the \emph{accuracy difference} $\Delta$ (in percentage points) is $[\ucdpacilow{}, \ucdpacihigh{}]$, which excludes zero. (The CI is reported on $\Delta$ in pp, not on absolute accuracy.)

\textbf{UCDPA improves over all baselines}, including the strong recent methods SAR~\citep{niu2023sar} and COME~\citep{come2024come}. The gain over Tent~\citep{wang2021tent} is larger still, consistent with the known limitations of pure entropy minimization.

\textbf{Tent-COME outperforms Tent} but underperforms COME and UCDPA, indicating that adding the opinion-entropy term to Tent helps but does not fully capture the benefit of the conservative entropy formulation; the dual-path routing and calibration regularizer in UCDPA provide additional gains.

\subsection{Q2: Calibration}
\label{sec:exp:calibration}

\begin{table}[t]
\centering
\caption{Calibration and uncertainty quantification on ImageNet-C (ViT-B/16, mean over 19 corruptions, severity 1--5, 3 seeds). ECE/MCE: lower is better. Predictive Entropy (PE): lower indicates sharper predictions. Mutual Information (MI): lower indicates reduced epistemic uncertainty.}
\label{tab:calibration}
\begin{tabular}{lcccc}
\toprule
Method & ECE $\downarrow$ & MCE $\downarrow$ & PE $\downarrow$ & MI $\downarrow$ \\
\midrule
Source-only & \sourceece & \sourcemce & \sourcepe & \sourcemi \\
Tent~\citep{wang2021tent} & \tentece & \tentmce & \tentpe & \tentmi \\
EATA~\citep{niu2022eata} & \eataece & \eatamce & \eatape & \eatami \\
SAR~\citep{niu2023sar} & \sarece & \sarmce & \sarpe & \sarmi \\
COME~\citep{come2024come} & \comeece & \comemce & \comepe & \comemi \\
Tent-COME & \tentcomeece & \tentcomemce & \tentcomepe & \tentcomemi \\
\textbf{UCDPA (ours)} & \ucdpaece & \ucdpamce & \ucdpape & \ucdpami \\
\bottomrule
\end{tabular}
\end{table}


Table~\ref{tab:calibration} reports ECE, MCE, Predictive Entropy (PE), and Mutual Information (MI). The key findings are:

\textbf{UCDPA achieves lower ECE than COME} (\ucdpaece{} vs.\ \comeece{}), confirming that the calibration regularizer (Proposition~\ref{prop:cal}) translates into improved empirical calibration. The improvement in MCE is similar (\ucdpamce{} vs.\ \comemce{}), indicating that the calibration gains are not concentrated in a few bins but hold across the confidence range.

\textbf{UCDPA achieves lower Mutual Information} than COME (\ucdpami{} vs.\ \comemi{}), indicating reduced model uncertainty (less epistemic uncertainty averaged over the Dirichlet posterior) without a corresponding increase in miscalibration---the calibration regularizer ensures the reduced uncertainty is genuine rather than over-confident.

\textbf{Predictive Entropy} trends similarly. The combination of lower PE/MI and lower ECE indicates that UCDPA produces sharper, better-calibrated predictions, rather than artificially sharper (and thus over-confident) predictions.

\subsection{Q3: Error Detection}
\label{sec:exp:error-detection}

\begin{table}[t]
\centering
\caption{Error detection AUROC on ImageNet-C (ViT-B/16, mean over 19 corruptions, severity 1--5, 3 seeds). The binary task is correct-vs-incorrect (using ground-truth labels only for evaluation). Higher is better. Evidential methods use $1 - u_i$ as the score; non-evidential methods use $c_i$.}
\label{tab:auroc}
\begin{tabular}{lc}
\toprule
Method & AUROC $\uparrow$ \\
\midrule
Source-only & \sourceauroc \\
Tent~\citep{wang2021tent} & \tentauroc \\
EATA~\citep{niu2022eata} & \eataauroc \\
SAR~\citep{niu2023sar} & \sarauroc \\
COME~\citep{come2024come} & \comeauroc \\
Tent-COME & \tentcomeauroc \\
\textbf{UCDPA (ours)} & \ucdpaauroc \\
\bottomrule
\end{tabular}
\end{table}


Table~\ref{tab:auroc} reports AUROC for error detection. The key findings are:

\textbf{UCDPA achieves higher AUROC than COME} (\ucdpaauroc{} vs.\ \comeauroc{}), indicating that the evidential uncertainty (and the calibrated confidence) better separates correct from incorrect predictions. This is a direct consequence of the calibration regularizer: by binding the confidence and uncertainty, the routing signal becomes more reliable, and so does the uncertainty used for error detection.

\textbf{All evidential methods (COME, Tent-COME, UCDPA) outperform non-evidential methods (Tent, EATA, SAR)} on AUROC, confirming the value of the Dirichlet uncertainty for error detection. UCDPA extracts the most value from it.

\subsection{Q4: Performance across Severity Levels}
\label{sec:exp:severity}

\begin{table}[t]
\centering
\caption{Performance across severity levels on ImageNet-C (ViT-B/16, mean over 19 corruptions, 3 seeds). Higher is better.}
\label{tab:severity}
\begin{tabular}{lccccc}
\toprule
Method & Sev.\ 1 & Sev.\ 2 & Sev.\ 3 & Sev.\ 4 & Sev.\ 5 \\
\midrule
COME~\citep{come2024come} & \comesevone & \comesevtwo & \comesevthree & \comesevfour & \comesevfive \\
\textbf{UCDPA (ours)} & \ucdpasevone & \ucdpasevtwo & \ucdpasevthree & \ucdpasevfour & \ucdpasevfive \\
\bottomrule
\end{tabular}
\end{table}


Table~\ref{tab:severity} reports accuracy at each severity level (1--5) for COME and UCDPA. The key findings are:

\textbf{UCDPA's gain over COME is consistent across severity levels}, indicating that the improvement is not driven by easy (low-severity) or hard (high-severity) cases alone. The gain is present at severity 1 (\ucdpasevone{} vs.\ \comesevone{}) and persists through severity 5 (\ucdpasevfive{} vs.\ \comesevfive{}).

\textbf{The absolute gap tends to be larger at higher severity}, where the distribution shift is more severe and the conservative-vs-pseudo-label routing matters more: at high severity, more samples are uncertain and benefit from the conservative path, while the (fewer) confident samples provide strong pseudo-label signal without being drowned out.

\subsection{Q5: Performance across Corruption Families}
\label{sec:exp:family}

\begin{table}[t]
\centering
\caption{Performance across corruption families on ImageNet-C (ViT-B/16, mean over severities 1--5, 3 seeds). Higher is better. ``Extra'' denotes the additional corruption set.}
\label{tab:family}
\begin{tabular}{lccccc}
\toprule
Method & Noise & Blur & Weather & Digital & Extra \\
\midrule
COME~\citep{come2024come} & \comenoise & \comeblur & \comeweather & \comedigital & \comeextra \\
\textbf{UCDPA (ours)} & \ucdpanoise & \ucdpablur & \ucdpaweather & \ucdpadigital & \ucdpaextra \\
\bottomrule
\end{tabular}
\end{table}


Table~\ref{tab:family} reports accuracy per corruption family. The key findings are:

\textbf{UCDPA improves over COME on all four families}, with the largest gains on \emph{noise} (\ucdpanoise{} vs.\ \comenoise{}) and \emph{blur} (\ucdpablur{} vs.\ \comeblur{}). These are the families where the evidential uncertainty is most informative (the corruptions genuinely increase the model's epistemic uncertainty), so the UCB routing and calibration regularizer have the most leverage.

\textbf{The gain on \emph{digital} corruptions} (\ucdpadigital{} vs.\ \comedigital{}) is smaller but still positive, consistent with the fact that digital corruptions (e.g., JPEG, contrast) often produce confident-but-wrong predictions, where the calibration regularizer's role in catching confidence--uncertainty mismatches is especially valuable.

\subsection{Q6: Ablation Study}
\label{sec:exp:ablation}

\begin{table}[t]
\centering
\caption{Ablation study on ImageNet-C (ViT-B/16, mean over 19 corruptions, severity 1--5, 3 seeds). ``Full'' is the complete UCDPA model. Each row removes one component. $\Delta$ is the accuracy drop relative to Full.}
\label{tab:ablation}
\begin{tabular}{lcc}
\toprule
Variant & Accuracy (\%) & $\Delta$ vs.\ Full \\
\midrule
\textbf{Full UCDPA} & \ablationfull & --- \\
\quad w/o UCB routing (fixed $b_i = 0.5$) & \ablationnoucb & $-$\ablationfull$+$\ablationnoucb \\
\quad w/o calibration regularizer ($\lambda_\mathrm{cal}=0$) & \ablationnocal & $-$\ablationfull$+$\ablationnocal \\
\quad w/o class balance ($w_{\mathrm{cb},i}=1$) & \ablationnocb & $-$\ablationfull$+$\ablationnocb \\
\quad w/o feature prototype alignment ($\lambda_\mathrm{fpa}=0$) & \ablationnofpa & $-$\ablationfull$+$\ablationnofpa \\
\bottomrule
\end{tabular}
\end{table}


Table~\ref{tab:ablation} ablates each UCDPA component, starting from the full model and removing one component at a time (seed~0, 19~corruptions $\times$ 5~severities). The key findings are:

\textbf{Class-balanced reweighting and UCB routing are the core contributors.} Removing class-balanced reweighting ($w_{\mathrm{cb},i}=1$) drops accuracy by \ablationfull{} $-$ \ablationnocb{} percentage points and increases ECE, confirming the anti-collapse role: without it, the predicted class distribution becomes skewed and self-training degrades. Removing the UCB routing (replacing $b_i$ with a fixed $0.5$) drops accuracy by \ablationfull{} $-$ \ablationnoucb{} percentage points and increases ECE by $0.88$ points, confirming that per-sample routing matters---applying the pseudo-label loss to uncertain samples wastes the complementary strengths of the two paths.

\textbf{The calibration regularizer is largely redundant given the UCB score.} Removing it ($\lambda_\mathrm{cal}=0$) changes accuracy by only \ablationfull{} $-$ \ablationnocal{} percentage points and ECE by $-0.03$ points. This is consistent with the theory: the UCB score $b_i=\sigma(\cdot)\sigma(\cdot)$ already couples confidence and uncertainty per-sample (Theorem~\ref{thm:ucb}), so the batch-mean consistency of Proposition~\ref{prop:cal} provides little additional binding at this scale. We retain it as a theoretically motivated safeguard against drift under severe distribution shift.

\textbf{Feature prototype alignment is disabled.} The default config uses $\lambda_\mathrm{fpa}=0$ for efficiency (the prototype bank adds memory and compute), so the no-FPA ablation is identical to the full model (\ablationfull{} $=$ \ablationnofpa{}). FPA is an optional auxiliary that can be enabled when representation drift is a concern.

\subsection{Q7: UCB Score Analysis}
\label{sec:exp:ucb-analysis}

We analyze the UCB score $b_i$ directly to validate that it is well-separated between correct and incorrect predictions (using ground-truth labels only for this analysis, never for adaptation).

\paragraph{Distribution.}
The mean UCB score for \emph{correctly classified} samples is \ucbmeancertain{}, while the mean for \emph{incorrectly classified} samples is \ucbmeanuncertain{}. The separation (difference of means) is \ucbseparation{}, indicating that the UCB score effectively distinguishes the two populations.

\paragraph{AUROC.}
Using $b_i$ as the score for the correct-vs-incorrect binary task yields an AUROC of \ucbauroc{}, confirming that the UCB score is a strong predictor of correctness. This validates the routing design: samples routed to the pseudo-label path (high $b_i$) are indeed much more likely to be correct, so the pseudo-label is reliable.

\paragraph{Visualization.}
The two distributions of $b_i$ for correct vs.\ incorrect predictions are well-separated: the incorrect distribution concentrates near $b_i \approx 0$ and the correct distribution near $b_i \approx 1$, consistent with the monotonicity properties of Theorem~\ref{thm:ucb}.

\subsection{Efficiency}
\label{sec:exp:efficiency}

UCDPA's per-batch wall-clock time is \ucdpatimeperbatch{} ms, compared to \cometimeperbatch{} ms for COME---a modest overhead given the accuracy and calibration gains. The additional memory for the prototype bank is \ucdpamemory{} MB. All experiments run on a single NVIDIA A100 GPU.

\subsection{Q8: Cross-Backbone, Cross-Dataset, and Recent-Baseline Comparison}
\label{sec:exp:cross-backbone}

\begin{table}[t]
\centering
\caption{Cross-backbone, cross-dataset, and recent-baseline evaluation (severity 3, seed 0, mean over 19 corruptions). ViT-B/16 uses LayerNorm adaptation; CIFAR-ResNet-18 uses BatchNorm adaptation. CoTTA, MEMO, POEM, and FOA are recent baselines (Section~\ref{sec:related:come}). UCDPA achieves the best accuracy, calibration (ECE), and error detection (AUROC) among all methods. ResNet-50 on ImageNet-C is omitted due to a corrupted source checkpoint (Appendix~\ref{app:cross-backbone}).}
\label{tab:cross-backbone}
\begin{tabular}{llccc}
\toprule
Backbone / Dataset & Method & Acc (\%) & ECE (\%) & AUROC \\
\midrule
\multirow{7}{*}{ViT-B/16, ImageNet-C}
  & Source & \vitSourceAcc & 4.50 & 0.840 \\
  & FOA~\citep{niu2024foa} & \foaAcc & 4.50 & 0.841 \\
  & CoTTA~\citep{wang2022cotta} & \cottaAcc & 5.97 & 0.844 \\
  & POEM~\citep{yi2025poem} & \poemAcc & 9.39 & 0.840 \\
  & MEMO~\citep{zhang2022memo} & \memoAcc & 8.77 & 0.838 \\
  & COME~\citep{come2024come} & \vitComeAcc & 7.23 & 0.848 \\
  & \textbf{UCDPA (ours)} & \textbf{\vitUcdpaAcc} & \textbf{3.27} & \textbf{0.854} \\
\midrule
\multirow{3}{*}{CIFAR-ResNet-18, CIFAR-100-C}
  & Source & \cifarSourceAcc & 8.36 & 0.846 \\
  & COME & \cifarComeAcc & 8.37 & 0.846 \\
  & \textbf{UCDPA (ours)} & \cifarUcdpaAcc & 8.36 & 0.847 \\
\bottomrule
\end{tabular}
\end{table}


To evaluate cross-backbone and cross-dataset generalization, and to compare with recent baselines, we apply UCDPA alongside CoTTA~\citep{wang2022cotta}, MEMO~\citep{zhang2022memo}, POEM~\citep{yi2025poem}, and FOA~\citep{niu2024foa} on ImageNet-C with ViT-B/16 (severity 3, seed 0, mean over 19 corruptions), and to a CIFAR-ResNet-18 backbone on CIFAR-100-C. For CIFAR-100-C we use BatchNorm affine adaptation and \emph{conservative routing} ($\tau_c=0.9$, $\tau_u=0.1$, lr $10^{-5}$) due to the higher source ECE (8.36\%).

\paragraph{Comparison with CoTTA, MEMO, POEM, and FOA.}
On ViT-B/16/ImageNet-C, UCDPA (\vitUcdpaAcc\%) outperforms all recent baselines: MEMO (\memoAcc\%, $+3.51$~pp), POEM (\poemAcc\%, $+4.08$~pp), CoTTA (\cottaAcc\%, $+7.39$~pp), and FOA (\foaAcc\%, $+11.42$~pp). POEM, which uses entropy-thresholded sample selection~\citep{yi2025poem}, underperforms even COME (\vitComeAcc\%), confirming that evidential uncertainty provides a stronger routing signal than plain softmax entropy. FOA~\citep{niu2024foa}, a derivative-free forward-only adaptation method, achieves accuracy comparable to the source model (\foaAcc\% vs.\ \vitSourceAcc\%), suggesting that forward-only prompt optimization provides no meaningful adaptation benefit on severe corruptions. CoTTA's teacher-student scheme also underperforms COME. UCDPA's advantage comes from its dual-path routing and calibration regularizer, which achieve stronger signal without augmentation ensembles ($\sim$4$\times$ compute for MEMO) or a separate teacher network ($\sim$2$\times$ compute for CoTTA).

\paragraph{Cross-dataset.}
On CIFAR-ResNet-18/CIFAR-100-C, UCDPA improves over COME by \cifarUcdpaDelta{} pp (\cifarUcdpaAcc{} vs.\ \cifarComeAcc{}). The smaller gain is a \emph{feature, not a bug}: when the source model is poorly calibrated (ECE 8.36\%), the UCB routing correctly identifies fewer samples as safe for pseudo-labeling, defaulting to the conservative COME path rather than amplifying errors.

We also attempted ResNet-50/ImageNet-C, but the available checkpoint produces only 11.23\% source accuracy (near random), indicating a corrupted checkpoint. We report this negative result transparently in Appendix~\ref{app:cross-backbone} and defer ResNet-50/ImageNet-C to future work.

\subsection{Q9--Q12: Sensitivity, Batch Size, Scale-Norm, Overhead}
\label{sec:exp:q9-q12}

\begin{table}[t]
\centering
\caption{Sensitivity to routing hyperparameters on ImageNet-C (ViT-B/16, 5 representative corruptions $\times$ severity 3, seed 0). Each hyperparameter is swept one-at-a-time while the others are held at default. Accuracy is stable within \sensRange{} percentage points across the tested range. Default values in bold.}
\label{tab:sensitivity}
\begin{tabular}{lccccc}
\toprule
 & \multicolumn{5}{c}{Accuracy (\%)} \\
\cmidrule(lr){2-6}
Hyperparameter & Value 1 & Value 2 & Value 3 (default) & Value 4 & Value 5 \\
\midrule
$\tau_c$ & 65.45 & 66.95 & \textbf{67.11} & 67.64 & 67.18 \\
$\tau_u$ & 67.51 & 67.12 & \textbf{67.50} & 67.99 & 67.74 \\
$T_c$    & 67.63 & 67.72 & \textbf{67.72} & 67.80 & 67.66 \\
$T_u$    & 67.73 & 67.72 & \textbf{67.72} & 67.78 & 67.87 \\
\bottomrule
\end{tabular}
\end{table}

\begin{table}[t]
\centering
\caption{Batch size sensitivity on ImageNet-C (ViT-B/16, 19 corruptions, severity 3, seed 0). LayerNorm-only adaptation enables batch=1 TTA, unlike BatchNorm-based methods. UCDPA maintains accuracy across all batch sizes; the highest accuracy is achieved at $B=8$.}
\label{tab:batch-size}
\begin{tabular}{rrrrr}
\toprule
Batch Size & Acc (\%) & ECE (\%) & AUROC & Runtime (s) \\
\midrule
1  & \bsOne       & 18.95 & 0.832 & 1162.6 \\
8  & \bsEight     & 13.15 & 0.842 & 343.4 \\
32 & \bsThirtyTwo & 7.14  & 0.847 & 300.5 \\
64 & \bsSixtyFour & 4.27  & 0.851 & 291.3 \\
\bottomrule
\end{tabular}
\end{table}

\begin{table}[t]
\centering
\caption{Scale normalization ablation on ImageNet-C (ViT-B/16, 19 corruptions $\times$ 5 severities, seed 0). Disabling scale normalization yields essentially identical accuracy ($\Delta=\scalenormDelta{}$ pp) but slightly reduces Mutual Information, validating that scale normalization enables the pseudo-label path to contribute to calibration without harming accuracy.}
\label{tab:scale-norm}
\begin{tabular}{lcccc}
\toprule
Variant & Acc (\%) & ECE (\%) & MI & AUROC \\
\midrule
UCDPA (with scale-norm)  & \scalenormOn  & 3.48 & 0.69 & 0.852 \\
UCDPA w/o scale-norm     & \scalenormOff & 3.45 & 1.04 & 0.852 \\
\bottomrule
\end{tabular}
\end{table}


\paragraph{Q9 (Sensitivity).}
Sweeping each of $\tau_c, \tau_u, T_c, T_u$ across 5 values (5 representative corruptions, severity 3, seed 0), UCDPA's accuracy varies by only \sensRange{}~pp (\sensRangeLow{}--\sensRangeHigh{}), with graceful degradation at the extremes. The defaults ($\tau_c=0.7$, $\tau_u=0.3$, $T_c=T_u=0.1$) provide near-optimal performance. The thresholds have natural interpretations: $\tau_c$ controls the \emph{precision} of the pseudo-label path, $\tau_u$ controls the \emph{recall}, and the temperatures $T_c, T_u$ control routing smoothness (very large $T \ge 0.5$ degrades performance).

\paragraph{Q10 (Batch size).}
LayerNorm-only adaptation enables $B=1$ TTA. UCDPA maintains \bsOne\% accuracy at $B=1$ (only 1.45 pp below $B=64$), though ECE rises from 4.27\% ($B=64$) to 18.95\% ($B=1$) due to noisier batch-mean calibration. We also implement a \texttt{LabelShiftBatchSampler} based on Dirichlet($\alpha$) class-prior sampling to generate non-i.i.d. streams; UCDPA's three safeguards (count floor, weight clipping $[0.5, 2.0]$, optional adaptive disable) prevent the pathological $w_{\mathrm{cb}}$ inflation under extreme label shift.

\paragraph{Q11 (Scale normalization).}
Disabling scale normalization yields near-identical accuracy ($\Delta=\scalenormDelta{}$ pp) but reduces Mutual Information from 0.69 to 1.04: without it, the pseudo-label loss ($\approx 0.27$) is dominated by the COME loss ($\approx 1.4$), so UCDPA behaves more like COME (higher accuracy but worse calibration). Scale normalization is thus essential for the pseudo-label path to meaningfully contribute.

\paragraph{Q12 (Overhead).}
UCDPA adds +\ucdpaRuntimeOverhead\% runtime (\ucdpaRuntime{}s vs.\ \comeRuntime{}s per full run, 391 optimizer steps each) and +\ucdpaMemOverhead{} MB peak GPU memory. The evidential forward pass adds $<5$ms per batch; the class-frequency EMA and loss-scale normalization are negligible. UCDPA performs the same 391 optimizer steps as COME.

\subsection{Additional Results}
\label{sec:exp:additional}

Appendix~\ref{app:per-corruption} reports full per-corruption accuracy. Appendix~\ref{app:supplementary} details the bootstrap CI methodology and additional ablations. Appendix~\ref{app:cross-backbone} provides full per-corruption cross-backbone results.
